% ===== PRÉAMBULE CANONIQUE — à coller VERBATIM en tête de CHAQUE .tex =====
\documentclass[11pt,a4paper]{article}
\usepackage[utf8]{inputenc}\usepackage[T1]{fontenc}\usepackage{lmodern}
\usepackage[french]{babel}
\usepackage{amsmath,amssymb,mathtools}\usepackage{icomma}
\usepackage{siunitx}\sisetup{locale=FR}  % \qty{}{}, \si{}, \ang{}
\usepackage{xcolor}\usepackage{colortbl}
\usepackage{tikz}\usepackage{pgfplots}\pgfplotsset{compat=1.18}
\usepackage[siunitx,europeanresistors]{circuitikz} % circuits (élec.)
\usepackage[most]{tcolorbox}\usepackage{enumitem}\usepackage{array,booktabs}
\usepackage[a4paper,margin=2.2cm]{geometry}
\usetikzlibrary{arrows.meta,calc,angles,quotes,positioning,patterns,%
  backgrounds,shapes.geometric,decorations.pathmorphing,decorations.markings,intersections}
\definecolor{primary}{RGB}{222,49,99}\definecolor{primarydark}{RGB}{150,22,55}
\definecolor{defcol}{RGB}{0,114,178}\definecolor{methcol}{RGB}{52,92,148}
\definecolor{excol}{RGB}{0,135,90}\definecolor{attncol}{RGB}{200,40,55}
\tcbset{boxbase/.style={enhanced,breakable,boxrule=0.9pt,arc=2.5pt,
  left=3mm,right=3mm,top=2.3mm,bottom=2.3mm,fonttitle=\bfseries,coltitle=white}}
% --- boîtes TUTORIEL (noms EXACTS, ne pas renommer) ---
\newtcolorbox{cle}[1][]{boxbase,colback=defcol!6,colframe=defcol,
  title={\textsc{À retenir}\ifx\relax#1\relax\relax\else\textmd{ : \itshape #1}\fi}}
\newtcolorbox{methode}[1][]{boxbase,colback=methcol!7,colframe=methcol,sharp corners,
  title={\textsc{Méthode}\ifx\relax#1\relax\relax\else\textmd{ --- \itshape #1}\fi}}
\newtcolorbox{exemplebox}[1][]{boxbase,colback=excol!5,colframe=excol,
  title={\textsc{Exemple}\ifx\relax#1\relax\relax\else\textmd{ : \itshape #1}\fi}}
\newtcolorbox{piege}[1][]{boxbase,colback=attncol!6,colframe=attncol,
  title={\textsc{Piège}\ifx\relax#1\relax\relax\else\textmd{ : \itshape #1}\fi}}
\newtcolorbox{remarque}[1][]{boxbase,colback=black!4,colframe=black!45,
  title={\textsc{Remarque}\ifx\relax#1\relax\relax\else\textmd{ : \itshape #1}\fi}}
% --- boîtes EXERCICE / CORRIGÉ (noms EXACTS, auto-comptées) ---
\newtcolorbox[auto counter]{exo}[1][]{boxbase,colback=primary!4,colframe=primary,
  title={\textsc{Exercice}~\thetcbcounter\ifx\relax#1\relax\relax\else\textmd{ --- \itshape #1}\fi}}
\newtcolorbox[auto counter]{corrige}[1][]{boxbase,colback=excol!4,colframe=excol,
  title={\textsc{Corrigé}~\thetcbcounter\ifx\relax#1\relax\relax\else\textmd{ --- \itshape #1}\fi}}
\newcommand{\vv}[1]{\vec{#1}}\newcommand{\ds}{\displaystyle}
\newcommand{\R}{\mathbb{R}}\newcommand{\N}{\mathbb{N}}\newcommand{\Z}{\mathbb{Z}}
% (un \usepackage{fancyhdr}+en-tête est possible ; il est ignoré côté web)
% =========================================================================
\begin{document}
\sisetup{per-mode=symbol}

\begin{corrige}[chute d'un objet lâché]
\begin{enumerate}[label=\alph*)]
  \item $h=\tfrac12 g t^2 \Rightarrow t=\sqrt{\dfrac{2h}{g}}=\sqrt{\dfrac{2\times20}{10}}=\sqrt{4}=\qty{2}{\second}$.
  \item $v=g\,t=10\times2=\qty{20}{\metre\per\second}$ (ou $v=\sqrt{2gh}=\sqrt{400}=\qty{20}{\metre\per\second}$).
\end{enumerate}
\end{corrige}

\begin{corrige}[vitesse et hauteur pendant la chute]
\begin{enumerate}[label=\alph*)]
  \item $v=g\,t=10\times3=\qty{30}{\metre\per\second}$.
  \item $h=\tfrac12 g t^2=\tfrac12\times10\times9=\qty{45}{\metre}$.
\end{enumerate}
\end{corrige}

\begin{corrige}[le temps de chute d'un tir horizontal]
\begin{enumerate}[label=\alph*)]
  \item Selon la verticale, c'est une chute libre ($v_{0y}=0$) : $t=\sqrt{\dfrac{2h}{g}}=\sqrt{\dfrac{2\times45}{10}}=\sqrt{9}=\qty{3}{\second}$.
  \item \textbf{Non} : le mouvement vertical (la chute) est indépendant du mouvement horizontal. La vitesse $v_{0x}$ n'agit que sur la portée, jamais sur la durée de chute.
\end{enumerate}
\end{corrige}

\begin{corrige}[portée d'un tir horizontal]
\begin{enumerate}[label=\alph*)]
  \item Horizontalement (MRU) : $x=v_{0x}\,t=15\times2=\qty{30}{\metre}$.
  \item Verticalement : $v_y=g\,t=10\times2=\qty{20}{\metre\per\second}$.
\end{enumerate}
\end{corrige}

\begin{corrige}[deux billes lâchées et lancées]
\begin{enumerate}[label=\alph*)]
  \item \textbf{Elles touchent le sol en même temps} : leur chute verticale est identique ($v_{0y}=0$ pour les deux), et $v_{0x}$ n'influe pas sur le temps de chute.
  \item La bille B (lancée horizontalement) : elle seule possède une vitesse horizontale, donc une portée non nulle. A tombe à la verticale.
\end{enumerate}
\end{corrige}

\end{document}
